\documentclass[11pt]{article}
\usepackage{restaurant-docs}

\begin{document}
\docmeta{API-documentatie}

\section{Gebruik, authenticatie en foutcodes}
Dit is de API van het restaurant. Lokaal draait hij op
\texttt{http://localhost:8080}. Importeer
\texttt{exports/restaurant.postman\_collection.json} en loop de requests per map
af.

Voor gewone logins en registratie: Authorization Code + PKCE (S256) via Keycloak
op \texttt{http://localhost:8081/realms/restaurant}. Client:
\texttt{restaurant-postman}; callback:
\texttt{https://oauth.pstmn.io/v1/callback}. Scope: \texttt{openid profile email}.

Lokale demoaccounts: \texttt{klant} / \texttt{Klant-Demo-2026!} (CUSTOMER) en
\texttt{personeel} / \texttt{Personeel-Demo-2026!} (STAFF). Nieuwe registraties
krijgen CUSTOMER. Dit zijn uitsluitend lokale demogegevens.

\begin{lstlisting}
Authorization: Bearer <access_token>
Content-Type: application/json
\end{lstlisting}

GET geeft 200; POST 201 met Location; PUT 200 (afbeelding 204); DELETE 204.
Fouten: 400 invoer, 401 token, 403 rechten/eigenaarschap, 404 ontbrekend, 405
methode, 409 conflict, 413 uploadgrootte, 415 bestandstype. Applicatiefouten
gebruiken ProblemDetail.

\begin{lstlisting}
{"type":"about:blank","title":"Conflict","status":409,
 "detail":"Table is reserved during this time slot"}
\end{lstlisting}

Ids zijn positieve database-identifiers. Datum/tijd is UTC zonder suffix;
bedragen zijn EUR-eindprijzen. Er zijn geen lijstfilters of paginering.

\section{Categorieën}
Categorie beheren. Verwijderen is alleen mogelijk zonder gekoppelde producten.
GET: CUSTOMER en STAFF. POST, PUT, DELETE: STAFF.

\texttt{GET /api/categories} --- lijst (200).
\texttt{GET /api/categories/\{id\}} --- één object (200/404).
\texttt{POST /api/categories} --- aanmaken (201 + Location).
\texttt{PUT /api/categories/\{id\}} --- wijzigen (200).
\texttt{DELETE /api/categories/\{id\}} --- verwijderen indien ongebruikt (204/409).

\begin{lstlisting}
POST {"name":"Desserts","description":"Nagerechten"}
201  {"id":3,"name":"Desserts","description":"Nagerechten"}
\end{lstlisting}

\section{Producten}
Producten beheren. Een gebruikt product kan niet worden verwijderd; zet
\texttt{available} op \texttt{false}. GET: CUSTOMER en STAFF. POST, PUT,
DELETE: STAFF.

\texttt{GET /api/products} --- lijst (200).
\texttt{GET /api/products/\{id\}} --- één product (200/404).
\texttt{POST /api/products} --- aanmaken (201 + Location).
\texttt{PUT /api/products/\{id\}} --- wijzigen (200).
\texttt{DELETE /api/products/\{id\}} --- verwijderen indien toegestaan (204/409).

\begin{lstlisting}
POST {"name":"Tiramisu","categoryId":1,"price":7.5,"available":true}
201  {"id":3,"name":"Tiramisu","categoryId":1,"price":7.5,
      "available":true,"imageUrl":"/api/products/3/image"}
\end{lstlisting}

\section{Tafels}
Tafels beheren. \texttt{available} geeft aan of een tafel beschikbaar is voor
gebruik. Reserveringen bepalen of een tafel op een bepaald tijdstip bezet is.
GET: CUSTOMER en STAFF. POST, PUT, DELETE: STAFF.

\texttt{GET /api/tables} --- lijst (200).
\texttt{GET /api/tables/\{id\}} --- één tafel (200/404).
\texttt{POST /api/tables} --- aanmaken (201 + Location).
\texttt{PUT /api/tables/\{id\}} --- wijzigen (200).
\texttt{DELETE /api/tables/\{id\}} --- verwijderen indien toegestaan (204/409).

\begin{lstlisting}
POST {"tableNumber":3,"seats":4,"available":true}
201  {"id":3,"tableNumber":3,"seats":4,"available":true}
\end{lstlisting}

\section{Reserveringen}
Klanten beheren eigen reserveringen; personeel mag alle reserveringen lezen,
wijzigen en annuleren. DELETE is een soft delete (annuleren). Duur: twee uur.

\texttt{GET /api/reservations} --- lijst (200).
\texttt{GET /api/reservations/\{id\}} --- één reservering (200/404).
\texttt{POST /api/reservations} --- aanmaken (201 + Location).
\texttt{PUT /api/reservations/\{id\}} --- wijzigen (200).
\texttt{DELETE /api/reservations/\{id\}} --- annuleren (204).

\begin{lstlisting}
POST {"customerName":"Ali","tableId":1,
      "startsAt":"2030-12-01T17:00:00","guests":2,
      "remarks":"Bij het raam"}
201  {"id":2,"customerName":"Ali","tableId":1,
      "startsAt":"2030-12-01T17:00:00",
      "endsAt":"2030-12-01T19:00:00","guests":2,
      "remarks":"Bij het raam","cancelled":false}
\end{lstlisting}

\section{Bestellingen}
Bestellingen beheren. Uitsluitend STAFF. Prijzen en productnamen worden
opgeslagen op het moment van bestellen. De server berekent het totaalbedrag.
Gefactureerde bestellingen kunnen niet meer worden gewijzigd.

\texttt{GET /api/orders} --- lijst (200).
\texttt{GET /api/orders/\{id\}} --- één bestelling (200/404).
\texttt{POST /api/orders} --- aanmaken (201 + Location).
\texttt{PUT /api/orders/\{id\}} --- wijzigen (200).
\texttt{DELETE /api/orders/\{id\}} --- verwijderen indien toegestaan (204/409).

\begin{lstlisting}
POST {"tableId":1,"lines":[{"productId":1,"quantity":2}]}
201  {"id":1,"tableId":1,"createdAt":"2030-12-01T17:10:00",
      "lines":[{"productId":1,"productName":"Pasta",
                "unitPrice":14.5,"quantity":2}],
      "total":29.0}
\end{lstlisting}

\section{Facturen}
Facturen beheren. Uitsluitend STAFF. Per bestelling kan maximaal één factuur
worden aangemaakt. Met PUT kan alleen de betaalstatus (\texttt{paid}) worden
gewijzigd. Betaalde facturen kunnen niet worden verwijderd of teruggezet naar
onbetaald.

\texttt{GET /api/invoices} --- lijst (200).
\texttt{GET /api/invoices/\{id\}} --- één factuur (200/404).
\texttt{POST /api/invoices} --- aanmaken (201 + Location).
\texttt{PUT /api/invoices/\{id\}} --- betaalstatus wijzigen (200).
\texttt{DELETE /api/invoices/\{id\}} --- onbetaalde factuur verwijderen (204/409).

\begin{lstlisting}
POST {"orderId":1}
201  {"id":1,"orderId":1,"createdAt":"2030-12-01T18:00:00",
      "total":29.0,"paid":false}
PUT  {"paid":true}
\end{lstlisting}

\section{Validatie per invoertype}
\begin{description}[style=nextline,leftmargin=2.4cm]
  \item[Categorie] \texttt{name} verplicht, niet leeg, max.\ 80 tekens, uniek;
        \texttt{description} verplicht, max.\ 500 tekens.
  \item[Product] \texttt{name} max.\ 100 tekens; \texttt{categoryId} bestaand;
        \texttt{price} niet negatief, max.\ acht gehele en twee decimale cijfers;
        \texttt{available} boolean.
  \item[Tafel] \texttt{tableNumber} $\ge 1$ en uniek; \texttt{seats} 1--50;
        capaciteitsverlaging geweigerd zolang reserveringen gekoppeld zijn.
  \item[Reservering] \texttt{customerName} max.\ 100 tekens; \texttt{startsAt}
        in de toekomst; \texttt{guests} 1--50 en passend bij de tafel;
        \texttt{remarks} max.\ 1000 tekens. Subject komt uit het token.
  \item[Bestelling] \texttt{lines} 1--100; product uniek en beschikbaar;
        \texttt{quantity} 1--100. \texttt{total} en \texttt{createdAt} zijn
        servervelden.
  \item[Factuur] POST: alleen \texttt{orderId} van een nog niet gefactureerde
        bestelling. PUT: boolean \texttt{paid}.
\end{description}

\section{Afbeeldingen en autorisatie}
\texttt{PUT /api/products/\{productId\}/image} --- STAFF, multipart veld
\texttt{file}, succes 204. Alleen PNG/JPEG tot 2~MiB en 16 miljoen pixels.

\begin{lstlisting}
curl.exe -X PUT http://localhost:8080/api/products/1/image ^
  -H "Authorization: Bearer TOKEN" -F "file=@product.png"
\end{lstlisting}

\texttt{GET /api/products/\{productId\}/image} --- CUSTOMER en STAFF, bytes met
\texttt{image/png} of \texttt{image/jpeg}. Ontbrekend: 404.

\begin{description}[style=nextline,leftmargin=3.2cm]
  \item[Niet ingelogd] GET \texttt{/} en \texttt{/api-docs.html}. API: 401.
  \item[CUSTOMER] Catalogus, tafels, afbeeldingen lezen; eigen reserveringen.
  \item[STAFF] Alle API-endpoints.
  \item[Onbekende rol] Geen API-toegang.
\end{description}

JWT's moeten geldig ondertekend zijn, de ingestelde issuer hebben en
\texttt{restaurant-api} in \texttt{aud} bevatten.

Bronnen: mijn ideefase, de eindopdracht, de Spring Security JWT-documentatie en
de Keycloak-import. OpenAI Codex heeft meegeholpen. De prompts staan in
\texttt{docs/ai-prompts.md}.

\end{document}
